\documentclass[conference]{IEEEtran}

\usepackage{amsmath,amssymb,amsfonts}
\usepackage{graphicx}
\usepackage{booktabs}
\usepackage{array}
\usepackage{multirow}
\usepackage{url}
\usepackage[hidelinks]{hyperref}
\usepackage{caption}
\usepackage{cite}

% --- Minor IEEE-ish tightening ---
\captionsetup{font=footnotesize}
\setlength{\textfloatsep}{10pt plus 2pt minus 2pt}
\setlength{\floatsep}{8pt plus 2pt minus 2pt}
\setlength{\intextsep}{8pt plus 2pt minus 2pt}

\title{Embodied Governance: Deterministic Governance Architecture with Bootstrap-Gated Escalation and Invariant Drift Detection}

\author{
\IEEEauthorblockN{AVC Systems Studio \quad / \quad Intuition Labs R\,+\,D}
\IEEEauthorblockA{Working Draft (IEEE-style)\\Timestamp: 02-28-2026 06:29 PM EST}
}

\begin{document}
\maketitle

\begin{abstract}
We formalize a deterministic governance architecture for safety-critical early-stage organizations operating under asymmetric capital and regulatory pressure.
The architecture integrates (i) a Structural Integrity Index (SSI), (ii) bootstrap confidence-interval (CI) gating for escalation decisions, (iii) a versioned invariant registry with drift detection, and (iv) an escalation finite-state machine (FSM) enforcing fail-closed defaults.
We position SSI as a governance coherence metric with simulation evidence and cross-domain plausibility; empirical instrument validation is treated as a single-domain next step.
\end{abstract}

\begin{IEEEkeywords}
governance systems, safety engineering, invariants, drift detection, bootstrap confidence intervals, finite-state machines
\end{IEEEkeywords}

% ============================================================
\section{Problem Definition}
Early-stage safety-technology organizations face a recurrent governance failure mode: constraints erode silently under investor urgency, operational load, or regulatory ambiguity.
Traditional governance artifacts (policies, committees, memos) are not state-aware and do not bind execution to invariants deterministically.
The result is escalation latency, authority instability, and un-audited exception-making.

This paper proposes an execution-oriented governance architecture: a small set of measurable invariants, a coherence metric (SSI), statistically gated escalation triggers, a deterministic FSM, and an invariant registry with drift detection.

% ============================================================
\section{Structural Integrity Index (SSI)}
\subsection{Definition}
Let $e_i \in [0,1]$ be the alignment score of event $i$ to its governing invariant(s), with $1$ meaning full compliance.
Let $w_i > 0$ be a criticality weight and $\Delta t_i \ge 0$ the age of the event.
Let $\kappa > 0$ be a recency decay constant.
Define $Z = \sum_{i=1}^{n} w_i$ and $\sigma(x) = \frac{1}{1+e^{-x}}$.

We define SSI at time $t$ as
\begin{equation}
\label{eq:ssi}
SSI_t \;=\; 10 \cdot \sigma\!\left(\frac{1}{Z}\sum_{i=1}^{n} w_i\, e_i\, e^{-\kappa \Delta t_i}\right).
\end{equation}

The logistic bound maps the normalized weighted alignment into $[0,10]$, enabling stable thresholding behavior and preventing unbounded amplification.

\subsection{Bootstrap CI Gating}
To avoid action on fragile estimates, we compute a nonparametric bootstrap CI for $SSI_t$.
For $B$ resamples (sampling with replacement), compute $SSI^{(b)}$ and define:
\begin{align}
CI_{\mathrm{lower}} &= \mathrm{percentile}(SSI^{(b)}, 2.5), \\
CI_{\mathrm{upper}} &= \mathrm{percentile}(SSI^{(b)}, 97.5).
\end{align}

\textbf{Gating rule (fail-closed):} a state transition that depends on SSI is permitted only if the conservative condition $CI_{\mathrm{lower}} < \tau$ is satisfied for an escalation threshold $\tau$ (e.g., $\tau=7.5$), i.e., we treat uncertainty as risk, not permission.

\subsection{Assumptions and Robustness Notes}
SSI is a coherence metric, not an empirical instrument by default.
The architecture assumes:
\begin{itemize}
\item \textbf{Windowing over stationarity:} SSI is computed over a rolling window of recent events rather than assuming global stationarity.
\item \textbf{Correlation awareness:} if event autocorrelation is detected, weights are penalized (e.g., $w'_i = w_i(1-\rho)$ for autocorrelation $\rho$ above a policy threshold).
\item \textbf{Auditability:} all event scores $e_i$, weight assignments $w_i$, and invariant versions are logged.
\end{itemize}

% ============================================================
\section{Escalation Finite-State Machine (FSM)}
We define a deterministic FSM with states:
\[
\texttt{IDLE} \rightarrow \texttt{ACTIVE} \rightarrow \texttt{STRESS} \rightarrow \texttt{FAULT} \rightarrow \texttt{RECOVERY}.
\]
A representative threshold policy is:
\begin{itemize}
\item $\texttt{ACTIVE} \rightarrow \texttt{STRESS}$ if $SSI_t < 7.5$.
\item $\texttt{STRESS} \rightarrow \texttt{FAULT}$ if $SSI_t < 5.0$ or quorum $M < 3$.
\item $\texttt{FAULT} \rightarrow \texttt{RECOVERY}$ only after invariant revalidation succeeds.
\end{itemize}

\textbf{Fail-closed default:} under invariant breach, unresolved drift, or insufficient validation confidence (e.g., CI ambiguity), execution is suspended or routed into protected states.

% ============================================================
\section{Regime Stability Result (Quorum Bound)}
We model witness quorum $M$ (independent validators) for escalation triggers.
Let $p_f \in (0,1)$ be the probability that a witness produces a spurious confirmation (false spike).

\textbf{Lemma 1 (False escalation bound).} Assuming independent witnesses,
\begin{equation}
\label{eq:pfm}
P(\text{false escalation}) = p_f^M.
\end{equation}

\textit{Sketch.} A false escalation requires all $M$ witnesses to fail simultaneously; independence implies multiplicative probability $p_f^M$.
Thus, increasing $M$ reduces false escalations exponentially.

This lemma supports conservative gating policies: quorum can be increased in high-stakes regimes to reduce false triggers without changing SSI itself.

% ============================================================
\section{Invariant Registry and Drift Detection}
SSI depends on the existence of a measurable invariant set.
To prevent SSI from collapsing into a subjective scoring function, invariants must be specified as structured, versioned objects with explicit tests and change control.

\subsection{Invariant Record}
Each invariant is stored as a signed versioned record:
\begin{itemize}
\item \textbf{Statement:} normative constraint (what must remain true).
\item \textbf{Scope:} systems/actions governed.
\item \textbf{Measurement:} function mapping logs/artifacts $\rightarrow e_i \in [0,1]$.
\item \textbf{Evidence policy:} required artifacts and provenance.
\item \textbf{Change control:} signatures, quorum, and veto policy.
\item \textbf{Review cadence / expiry:} scheduled revalidation triggers.
\end{itemize}

\subsection{Drift Detection}
Drift detection monitors (i) rolling alignment distribution shifts, (ii) weight instability, (iii) autocorrelation spikes, and (iv) unauthorized changes.
A breach triggers a deterministic FSM transition into \texttt{STRESS}.
Drift detection is conceptually distinct from low alignment: drift indicates the \emph{meaning} or \emph{measurement} of compliance is changing, which must be fail-closed by default.

% ============================================================
\section{Simulation Summary (Positioned as Plausibility Evidence)}
A controlled Monte Carlo evaluation (10,000 iterations) across multiple governance environments indicates $\sim$29--30\% reduction in failure metrics under conservative parameterization.
We emphasize: such outcomes demonstrate portability in simulation, not multi-domain empirical validation.

\begin{table}[t]
\caption{Illustrative Monte Carlo Summary (10,000 iterations)}
\label{tab:mc}
\centering
\begin{tabular}{lccc}
\toprule
Domain & Burnout Red. & Drift Red. & Custody Gain \\
\midrule
AI Governance & 29.9\% & 30.1\% & 29.8\% \\
Biotech Governance & 30.0\% & 29.7\% & 30.2\% \\
Blockchain Governance & 29.8\% & 30.0\% & 29.9\% \\
\bottomrule
\end{tabular}
\end{table}

% ============================================================
\section{Figures}
Fig.~\ref{fig:ssi} illustrates the logistic bounding curve used in Eq.~\ref{eq:ssi}.
Fig.~\ref{fig:fsm} shows threshold regions for the representative FSM policy.
Fig.~\ref{fig:drift} provides an illustrative drift visualization.

\begin{figure}[t]
\centering
\includegraphics[width=\columnwidth]{figures/ssi_curve.pdf}
\caption{SSI logistic bounding function.}
\label{fig:ssi}
\end{figure}

\begin{figure}[t]
\centering
\includegraphics[width=\columnwidth]{figures/fsm_thresholds.pdf}
\caption{FSM threshold regions (representative).}
\label{fig:fsm}
\end{figure}

\begin{figure}[t]
\centering
\includegraphics[width=\columnwidth]{figures/invariant_drift.pdf}
\caption{Invariant drift illustration via rolling alignment distributions (illustrative).}
\label{fig:drift}
\end{figure}

% ============================================================
\section{Conclusion}
We presented an execution-oriented governance architecture integrating SSI, bootstrap CI gating, a drift-aware invariant registry, and deterministic FSM enforcement.
SSI is framed as a coherence metric with simulation support; the publication-critical next step is anchored empirical validation in a single domain (e.g., AI governance firms), with cross-domain results treated as controlled transfer demonstrations.

% ============================================================
\appendices
\section{SSI and Bootstrap Pseudocode (Reference)}
\begin{verbatim}
logistic(x):
  return 1 / (1 + exp(-x))

calculate_ssi(events e, weights w, kappa, delta_t):
  Z = sum(w)
  S = sum_i w_i * e_i * exp(-kappa * delta_t_i)
  return 10 * logistic(S / Z)

bootstrap_ci_ssi(e, w, kappa, delta_t, B=1000):
  for b in 1..B:
    sample indices with replacement
    SSI^(b) = calculate_ssi(sample)
  CI_lower = percentile(2.5)
  CI_upper = percentile(97.5)
\end{verbatim}

\bibliographystyle{IEEEtran}
\bibliography{refs}

\end{document}
